\begin{table}[H]
\centering
\caption{Standardized Effect Sizes for Main Outcomes}
\label{tab:sde}
\small
\begin{threeparttable}
\begin{tabular}{@{}lrrcrrcl@{}}
\toprule
Outcome & $\hat{\beta}$ & SE & SD($X$) & SD($Y$) & SDE & SE(SDE) & Classification \\
\midrule
FP Emp (Ban) & $-$0.008 & 0.161 & --- & 1.312 & $-$0.006 & 0.123 & Small neg. \\
FP Emp (Recv) & 0.254 & 0.138 & --- & 1.113 & 0.228 & 0.124 & Large pos. \\
Phys Emp (Ban) & $-$0.031 & 0.026 & --- & 1.087 & $-$0.028 & 0.024 & Small neg. \\
FP Emp (DDD) & $-$0.055 & 0.154 & --- & 2.275 & $-$0.024 & 0.068 & Small neg. \\
\bottomrule
\end{tabular}
\begin{tablenotes}[flushleft]
\scriptsize
\item \textit{Notes:}
\textbf{Country:} United States.
\textbf{Research question:} Did state-level abortion bans following the Dobbs decision
cause reallocation of reproductive healthcare employment from ban states to neighboring
receiving states?
\textbf{Policy mechanism:} Thirteen states had pre-enacted trigger laws that immediately
banned most abortions upon the Supreme Court's reversal of Roe v.\ Wade in June 2022,
causing family planning clinics to close or reduce services in ban states while demand
concentrated in bordering non-ban states.
\textbf{Outcome definition:} Log beginning-of-quarter employment from Census QWI
for family planning centers (NAICS 62141) and physician offices (NAICS 6211).
\textbf{Treatment:} Binary indicator for state abortion ban in effect (ban states)
or bordering a ban state (receiving states).
\textbf{Data:} Census QWI, 2015Q1--2024Q2, state-quarter level.
\textbf{Method:} TWFE DiD with state and quarter FE, state-clustered SEs.
\textbf{Sample:} All US states; partial-ban states (GA, IN, SC, OH) in control group.
SDE $= \hat{\beta} / \text{SD}(Y)$ where SD($Y$) is pre-treatment SD.
Classification refers to magnitude, not statistical significance:
Large ($|$SDE$| > 0.15$), Moderate ($0.05$--$0.15$), Small ($0.005$--$0.05$),
Null ($< 0.005$).
\end{tablenotes}
\end{threeparttable}
\end{table}
